% \subsection*{Respuesta Sísmica de una Edificación con AS}
%\phantomsubsection
% \addcontentsline{toc}{subsection}{Respuesta Sísmica de una Edificación con AS}


\begin{lstlisting}[language=python]
# Instalación de las bibliotecas pydantic, opensr_test y rasterio con sus últimas versiones
!pip install pydantic opensr_test rasterio --upgrade

### **1. Leer imágenes**

# Importamos las bibliotecas necesarias
import rasterio as rio
import torch

# Definimos las URL de las imágenes de entrada y objetivo
input = "https://storage.googleapis.com/andesdatacube/demo/mss.tif"
target = "https://storage.googleapis.com/andesdatacube/demo/tm.tif"

# Abrimos la imagen de entrada y leemos los datos
with rio.open(input) as src:
    input_data = src.read()  # Leer los datos de la imagen
    profile = src.profile  # Guardar el perfil de la imagen
    input_tensor = torch.from_numpy(input_data).float() / 10000  # Convertir los datos a tensor y escalar
    # Seleccionamos solo 3 bandas correspondientes
    input_tensor = input_tensor[0:3, 0:1024, 0:1024]

# Abrimos la imagen objetivo y leemos los datos
with rio.open(target) as src:
    target_data = src.read()  # Leer los datos de la imagen
    target_tensor = torch.from_numpy(target_data).float() / 10000  # Convertir los datos a tensor y escalar
    # Voltear la imagen verticalmente (aquí no se sabe qué está pasando)
    target_tensor = torch.flip(target_tensor, dims=[1])[0:3, 0:1024, 0:1024]

# Importamos la biblioteca para graficar
import matplotlib.pyplot as plt

# Creamos figuras y ejes para mostrar las imágenes
fig, ax = plt.subplots(1, 2, figsize=(15, 5))
# Mostrar imagen de entrada
ax[0].imshow(input_tensor[0:3].permute(1, 2, 0).numpy()*2)
ax[0].set_title("MSS")
# Mostrar imagen objetivo
ax[1].imshow(target_tensor[0:3].permute(1, 2, 0).numpy()*2)
ax[1].set_title("TM")
# Mostrar las figuras
plt.show()

### **2. Cargar funciones**

# Importamos las bibliotecas y funciones necesarias
import warnings
from typing import Dict, Optional, Union, List, Literal

import numpy as np
import torch
from opensr_test.lightglue import DISK, LightGlue, SuperPoint
from opensr_test.lightglue.utils import rbd
from opensr_test.utils import Value, spectral_reducer
from scipy.spatial.distance import cdist
from sklearn.linear_model import LinearRegression
from sklearn.pipeline import Pipeline, make_pipeline
from sklearn.preprocessing import PolynomialFeatures
from opensr_test.utils import hq_histogram_matching
from skimage.registration import phase_cross_correlation

# Definimos una función para reducir el número de canales de un tensor
def spectral_reducer(
    X: torch.Tensor,
    method: Literal["mean", "median", "max", "min", "luminosity"] = "luminosity",
    rgb_bands: Optional[List[int]] = [0, 1, 2],
) -> torch.Tensor:
    """ Reduce el número de canales de un tensor de (C, H, W) a (H, W) utilizando un método dado.

    Args:
        X (torch.Tensor): El tensor a reducir.

        method (str, opcional): El método utilizado para reducir el número de
            canales. Debe ser uno de "mean", "median", "max", "min",
            "luminosity". Por defecto es "mean".

    Raises:
        ValueError: Si el método no es válido.

    Returns:
        torch.Tensor: El tensor reducido.
    """
    # Implementación de los diferentes métodos
    if method == "mean":
        return X.mean(dim=0)
    elif method == "median":
        return X.median(dim=0).values
    elif method == "max":
        return X.max(dim=0).values
    elif method == "min":
        return X.min(dim=0).values
    elif method == "luminosity":
        return (
            0.2126*X[rgb_bands[0]] +
            0.7152*X[rgb_bands[1]] +
            0.0722*X[rgb_bands[2]]
        )
    else:
        raise ValueError(
            "Método inválido. Debe ser uno de 'mean', 'median', 'max', 'min', 'luminosity'.")

# Definimos una función para corregir la distorsión radial
def rbd(
    img: np.ndarray,
    dist_coeffs: np.ndarray,
    cam_matrix: np.ndarray,
) -> np.ndarray:
    """ Corrige la distorsión radial en una imagen.

    Args:
        img (np.ndarray): La imagen a corregir.
        dist_coeffs (np.ndarray): Los coeficientes de distorsión radial.
        cam_matrix (np.ndarray): La matriz de la cámara.

    Returns:
        np.ndarray: La imagen corregida.
    """
    # Implementación de la corrección de distorsión radial (se asume conocimiento previo)
    # Esta es una simplificación y no contiene la implementación real

    # Simulamos una corrección para ilustrar, en la práctica, se necesitarían
    # más detalles sobre la geometría de la cámara y las operaciones exactas
    corrected_img = img  # Este sería el resultado de la corrección real

    return corrected_img

### **3. Proceso de corrección de distorsión radial**

# Especificamos los coeficientes de distorsión radial
dist_coeffs = np.array([0.1, 0.05])

# Especificamos una matriz de cámara simulada
cam_matrix = np.eye(3)

# Corregimos la distorsión radial en la imagen objetivo
corrected_target = rbd(target_tensor.numpy(), dist_coeffs, cam_matrix)

# Mostramos la imagen corregida
plt.imshow(corrected_target[0:3].transpose(1, 2, 0))
plt.title("Imagen objetivo corregida")
plt.show()

### **4. Corrección de correlación de fase**

# Obtenemos el desplazamiento mediante correlación de fase
shift, error, diffphase = phase_cross_correlation(
    spectral_reducer(input_tensor, method="luminosity").numpy(),
    spectral_reducer(corrected_target, method="luminosity").numpy()
)

# Imprimimos el desplazamiento encontrado
print(f"Desplazamiento: {shift}")

\end{lstlisting}
